\begin{textsample}{POS Dim 3 – rn – Score 21.00 – rn/rn\_inf\_28.txt}  \label{ex:f3_pos_053}
7.
ORGANIZAÇÃO DE TEMPOS, ESPAÇOS, MATERIAIS E RELAÇÕES > Nossas \textbf{crianças} têm direito à \textbf{brincadeira} ; Nossas \textbf{crianças} têm direito à atenção individual ; Nossas \textbf{crianças} têm direito a um ambiente aconchegante, seguro e estimulante ; Nossas \textbf{crianças} têm direito ao contato com a natureza ; [... ] Nossas \textbf{crianças} têm direito a desenvolver sua \textbf{curiosidade}, imaginação e capacidade de expressão ; Nossas \textbf{crianças} têm direito ao movimento em espaços amplos ; Nossas \textbf{crianças} têm direito à proteção, ao \textbf{afeto} e à amizade ; Nossas \textbf{crianças} têm direito a expressar seus sentimentos ; [... ] ( CAMPOS, 2009, p. 13 ).
Objetiva -se neste capítulo, apresentar alguns apontamentos referentes à organização de tempos, espaços, materiais e relações, que podem constituir práticas pedagógicas na Educação \textbf{Infantil} no contexto da discussão atual acerca de currículos, as quais se materializam nas experiências concretas da vida cotidiana de \textbf{bebês}, \textbf{crianças} bem \textbf{pequenas} e \textbf{crianças} \textbf{pequenas}, tendo como eixos norteadores as interações e a \textbf{brincadeira}.
Nesta visão de currículo, a transmissão e a reelaboração de conhecimentos socialmente significativos estão presentes nos modos como se organizam as práticas pedagógicas, considerando as singularidades das \textbf{crianças}, bem como as diversidades culturais, sociais, etárias e políticas de cada grupo : > Tal definição reconhece tanto a intencionalidade da prática pedagógica na organização dos espaços, tempos, materiais, relações sociais, na seleção de experiências de aprendizagem quanto, e principalmente, no protagonismo de cada e toda \textbf{criança} nessa organização e seleção, em vez de ser um processo de ensino centrado em decisões apenas dos \textbf{adultos} ( BARBOSA ; OLIVEIRA, 2016, p. 23 ).
Considerando essa perspectiva, assume -se nesta proposta, o desafio de redimensionar as ações de planejamento docente buscando dialogar com o protagonismo das \textbf{crianças} na diversidade de experiências, interações, \textbf{brincadeiras}, linguagens e conhecimentos as quais podem constituir as práticas cotidianas na Educação \textbf{Infantil}.
Assim, > o planejamento curricular deve assegurar condições para a organização do tempo cotidiano das instituições de Educação \textbf{Infantil} de modo a equilibrar continuidade e inovação nas atividades, \textbf{movimentação} e concentração das \textbf{crianças}, momentos de segurança e momentos de desafio na participação das mesmas, e articular seus ritmos individuais, vivências pessoais e experiências coletivas com \textbf{crianças} e \textbf{adultos} ( BRASIL, 2009, p. 12 ).
Assumem -se aqui, alguns aspectos a serem considerados na organização das experiências a serem realizadas no cotidiano das unidades de Educação \textbf{Infantil}, de modo que este se constitua como contexto de vivência, aprendizagem e desenvolvimento das \textbf{crianças} : - os tempos de realização das atividades ( ocasião, frequência, duração ) ; - os espaços em que essas atividades transcorrem ( o que inclui a estruturação dos espaços internos, externos, de modo a favorecer as interações \textbf{infantis} na exploração que fazem do mundo ) ; - os materiais disponíveis ; - e, em especial, as maneiras de o professor exercer seu papel ( organizando o ambiente, ouvindo as \textbf{crianças}, respondendo -lhes de determinada maneira, oferecendo -lhes materiais, sugestões, apoio emocional, ou promovendo condições para a ocorrência de valiosas interações e \textbf{brincadeiras} criadas pelas \textbf{crianças} etc. ) ( OLIVEIRA, 2014, p. 193 ).
Para tanto, as \textbf{rotinas} pedagógicas que se estruturam diariamente nas instituições de Educação \textbf{Infantil} devem considerar a alternância de atividades livres e dirigidas, individuais e coletivas, garantindo experiências de aprendizagem em que o professor pode mediar e \textbf{apoiar} as interações e as atividades \textbf{lúdicas} das \textbf{crianças}.
Sugere -se então : - transformar o espaço físico em um ambiente acolhedor, estimulante e desafiador, propício às investigações \textbf{infantis} ; - observar que os tempos das aprendizagens das \textbf{crianças} \textbf{pequenas} exigem lentidão, continuidade, duração e também leveza, que não podem ser rápidos e fragmentados como o tempo da produção industrial, pois as \textbf{crianças} \textbf{pequenas} precisam de tempo para construir suas histórias e identidades pessoais e sociais ; - oferecer materialidades diversas, pois conhecimentos e indagações que emergem dos objetos e das relações possíveis, que se estabelecem desde um ambiente rico em oportunidades de interações materiais e imateriais, certamente favorecerão questionamentos e conhecimentos... ( BARBOSA ; OLIVEIRA, 2016, p. 27 ).
O desafio que se coloca para a elaboração curricular, consiste em transcender a prática pedagógica centrada no(a) professor(a) e garantir experiências que possibilitem “ o encontro de explicações pela \textbf{criança} sobre o que ocorre à sua \textbf{volta} e consigo mesma, enquanto desenvolve formas de sentir, pensar e solucionar problemas ” ( OLIVEIRA, 2013, p. 06 ).
Sendo assim, os planejamentos das experiências precisam assegurar um caráter vivencial, interdisciplinar e mobilizador da autonomia \textbf{infantil}.
Ser produzidos para uma realização \textbf{diária}, semanal, mensal ou por períodos mais longos, no caso de projetos.
Ser acompanhados e avaliados pelo \textbf{registro} continuado das observações críticas e criativas dos professores sobre as atividades realizadas, em especial das \textbf{brincadeiras} e interações das \textbf{crianças} ( OLIVEIRA, 2015, p. 22 ).
Propõe -se então, que o planejamento seja resultado da observação e documentação que o(a) professor(a) constrói a partir da relação cotidiana com as \textbf{crianças}.
Sobre isso, orienta -se que : > No seu planejamento \textbf{diário}, o professor deve destinar tempo às propostas que ele fará ao grupo e também tempo para que as próprias \textbf{crianças} inventem seus problemas, coloquem -se desafios de criação, desenvolvam seus projetos pessoais em qualquer situação, seja em ateliês de arte, no parque etc.
A \textbf{criança} tem muito que aprender sobre seu tempo de produção e o professor, consequentemente, deve organizar modos de \textbf{apoiar} essa importante aprendizagem.
Observando atentamente como as \textbf{crianças} vivem o tempo de criação, será possível criar alternativas à gestão da sala, para que não haja homogeneização desnecessária ( AUGUSTO, 2013, p. 27 ).
Salles e Faria ( 2012 ), assim como, Oliveira et al ( 2012 ), propõem a gestão do tempo curricular tendo como base, modalidades organizativas.
Desse modo, orientam que pode acontecer a instituição de tempos mais estáveis e permanentes para as atividades, o que propõe uma aproximação com vistas à construção de familiaridade com determinadas práticas, que exigem o desenvolvimento de \textbf{hábitos} e comportamentos específicos.
É o caso, por exemplo, da roda para \textbf{conversar}, da roda para ler e \textbf{contar} histórias e dos momentos de alimentação.
Assim como, a proposição de sequências didaticamente pensadas para inserir graus crescentes de desafios às \textbf{crianças}, sempre baseadas em avaliações das aprendizagens e na projeção de novos objetivos.
E, a organização de projetos coletivos que permitem à \textbf{criança} aprender com seus pares e ser \textbf{apoiada} na pesquisa, investigação, sistematização e comunicação de novos conhecimentos, utilizando seus próprios recursos, além de outros que ela pode ter acesso no ambiente da Educação \textbf{Infantil}, tais como livros, vídeos, instrumentos e materiais específicos, etc.
Estas modalidades metodológicas envolvem, ou podem envolver, conjuntos de experiências educativas.
São modalidades desenvolvidas no contexto da prática pedagógica que funcionam como mecanismos organizativos da gestão do tempo das \textbf{crianças} na instituição.
As atividades-experiências permanentes são tempos mais estáveis e constantes que ocorrem com regularidade ( \textbf{diária}, semanal ou quinzenal ) e têm como objetivo familiarizar as \textbf{crianças} com determinadas experiências de aprendizagem, que promovem o desenvolvimento de habilidades, procedimentos e comportamentos relativos aos processos de apropriação de linguagens e construção de conhecimentos, que demandam constância, diversidade e continuidade em sua realização.
São desenvolvidas de forma regular, sistemática e previsível, e precisam acontecer em espaços-tempos previamente organizados e propostos periodicamente.
Podem propor aos educandos uma familiarização com diversas práticas-experiências como maneiras de ler, de \textbf{brincar}, de produzir textos, de fazer artes, de experimentar, de argumentar, etc.
As sequências didáticas, por sua vez, se constituem como um conjunto de experiências, didaticamente sequenciadas e articuladas, considerando sua continuidade, diversidade e contextualização para sistematizar linguagens e conhecimentos específicos, cujas aprendizagens requerem aprimoramento com a experiência.
Ou seja, o que acontece e se desenvolve a seguir, depende do que já foi realizado e aprendido.
São, no entanto, os projetos que precisam ser especialmente desenvolvidos com as \textbf{crianças}.
Referem -se a um conjunto de experiências realizadas em torno de um objetivo compartilhado com elas a partir de suas questões acerca da temática em investigação.
Tal modalidade metodológica promove a participação das \textbf{crianças} como sujeitos ativos e investigativos, oportunizando processos de apropriação e construção de conhecimentos mediados pelos \textbf{adultos} e pela linguagem, em contextos significativos.
Os projetos têm, portanto, como fios estruturantes, a pesquisa, os \textbf{registros} sistemáticos em diferentes linguagens e a socialização.
As experiências desenvolvidas durante um projeto de pesquisa com as \textbf{crianças} devem estar encadeadas em função da problematização e de sua constante retomada no processo.
Assim sendo, devem promover o intercâmbio de ideias, saberes e conhecimentos com vistas à sua ampliação e avaliação coletiva.
Para Katz ( 1999 ), os projetos oferecem a parte do currículo na qual as \textbf{crianças} são encorajadas a tomarem suas próprias decisões e a fazerem suas próprias escolhas, geralmente em cooperação com seus \textbf{colegas}, sobre o trabalho a ser realizado. > Este tipo de trabalho aumenta a confiança das \textbf{crianças} em seus próprios poderes intelectuais e reforça sua disposição de continuar aprendendo.
As atividades incluem observação direta, perguntas a pessoas e a especialistas relevantes, coleta de artefatos pertinentes, representação de observações, de ideias, de memórias, de emoções, de imagens e de novos conhecimentos em várias maneiras, incluindo encenação \textbf{dramática} ( KATZ, 1999, p. 38 ).
Trata -se de um percurso coletivo e articulado, em que as \textbf{crianças} tornam -se autoras, coautoras e corresponsáveis pela própria aprendizagem.
O trabalho com projetos exige práticas de \textbf{registros} diversos, criando espaços e instrumentos de ação, reflexão e investigação da prática educativa, sendo o educador aprendiz e protagonista em parceria com o grupo com o qual atua, assim como sujeito ativo dos processos de experimentações e descobertas.

% matched lemmas: adulto, afeto, apoiar, bebê, brincadeira, brincar, colega, contar, conversar, criança, curiosidade, diário, dramático, hábito, infantil, lúdico, movimentação, pequeno, registro, rotina, volta
\end{textsample}
